
% /science/math.tex

\chapter{Interlude on the mathematical sciences}

This interlude reviews the mechanics developed in the previous Chapter \ref{chap:mechanics} and introduces a couple of topics from the mathematical sciences needed for Chapter \ref{chap:quantum}.

By Newton's second law of motion, the net force $F$ acting on a body with constant mass $m \geq 0$ and mechanical momentum $p = m \dot{x}$ satisfies $F = \frac{dp}{dt}$.
If $m > 0$, then $\frac{dp}{dt} = m \frac{d}{dt}\p{\dot{x}}$ and Newton's second law yields $F = m \ddot{x}$.
If $m = 0$, then the mechanical momentum is $p = 0$ and the net force on the body is $F = 0$; Newton's second law then fails to determine the motion of the body.
For Newton's second law to determine the motion of a body, its mass must be strictly positive.

\section{Units}\label{sec:units}

Let $\vq = \begin{pmatrix} q_1 \\ \vdots \\ q_i \\ \vdots \\ q_n \end{pmatrix}$ denote a vector of $n$ generalized coordinates and $\vp = \begin{pmatrix} p_1 \\ \vdots \\ p_i \\ \vdots \\ p_n \end{pmatrix}$ denote their canonical momenta with respect to some Lagrangian $L$.

\begin{remark}[Units of derivatives of function on phase space]\label{rem:derivative}
	For any function $F\p{\vq, \vp}$ of generalized coordinates $\vq$ and canonical momenta $\vp = \frac{\partial L}{\partial \dot{\vq}}$, the partial derivative $\frac{\partial F\p{\vq, \vp}}{\partial q_i}$ with respect to any generalized coordinate $q_i$ has units $\left|\frac{\partial F\p{\vq, \vp}}{\partial q_i}\right| = \frac{\left|F\right|}{\left|q_i\right|}$ and the partial derivative $\frac{\partial F\p{\vq, \vp}}{\partial p_i}$ with respect to the canonical momentum $p_i$ conjugate to $q_i$ has units $\left|\frac{\partial F\p{\vq, \vp}}{\partial p_i}\right| = \frac{\left|F\right|}{\left|p_i\right|}$.
\end{remark}

\begin{proposition}\label{prop:action}
	For any generalized coordinate $q$ and canonical momentum $p = \frac{\partial L}{\partial \dot{q}}$, the units of the product $q p$ are $\left|q p\right| = \left|L\right| \left|t\right|$.
\end{proposition}
\begin{proof}
	The result follows from Remark \ref{rem:derivative}.
\end{proof}

\begin{lemma}\label{lem:derivatives}
	For any functions $F\left(\vq, \vp\right)$ and $G\left(\vq, \vp\right)$, the products of partial derivatives have the dimensions below.
	\begin{align*}
		\left|\frac{\partial F}{\partial q_i} \frac{\partial G}{\partial p_i}\right| & = \frac{\left|F\right| \left|G\right|}{\left|L\right| \left|t\right|}
		                                                                             &
		\left|\frac{\partial G}{\partial q_i} \frac{\partial F}{\partial p_i}\right| & = \frac{\left|G\right| \left|F\right|}{\left|L\right| \left|t\right|}
	\end{align*}
\end{lemma}
\begin{proof}
	The result follows from Remark \ref{rem:derivative} and Proposition \ref{prop:action}.
\end{proof}

The two products have the same dimensions.

\begin{theorem}\label{thm:Poisson}
	For any functions $F\left(\vq, \vp\right)$ and $G\left(\vq, \vp\right)$, the dimensions of the Poisson bracket $\left\{F, G\right\}$ are $\left|\{F,G\}\right| = \frac{\left|F\right| \left|G\right|}{\left|L\right| \left|t\right|}$.
\end{theorem}
\begin{proof}
	The result follows from Lemma \ref{lem:derivatives}.
\end{proof}

The units of the reduced Planck constant $\hbar$ are $\left|\hbar\right| = \left|L\right| \left|t\right|$.

\begin{corollary}\label{corl:Poisson}
	For any functions $F\left(\vq, \vp\right)$ and $G\left(\vq, \vp\right)$, the dimensions of $i \hbar \left\{F, G\right\}$ are $\left|i \hbar \{F,G\}\right| = \left|F\right| \left|G\right|$.
\end{corollary}
\begin{proof}
	The result follows from Theorem \ref{thm:Poisson}.
\end{proof}

\section{Algebra}

Given the axis spanned by any unit vector $\hat{n}$, the rotation $r_{\hat{n}}^\theta$ of $\R^3$ about the axis by $\theta$ can be expressed by Rodrigues's rotation formula below.
\begin{align}\label{eqn:Rodrigues}
	r_{\hat{n}}^\theta \vec{m} & = \vec{m} \cos \theta + \left(\hat{n} \times \vec{m} \right) \sin \theta + \hat{n} \left(\hat{n} \cdot \vec{m}\right)\left(1 - \cos \theta\right)
\end{align}

Let $\mA$ denote an operator on a finite-dimensional vector space.

Recall that the \emph{exponential} of $\mA$ is the operator $\exp\p{\mA} = e^{\mA}$ defined by $\dsp \sum_{k=0}^\infty \frac{1}{k!} \mA^k$.

\subsubsection*{Linear algebra}

Recall \emph{involutivity} of $\mA$ is the condition that $\mA^2 = I$.
If $\mA$ is an involution, then it holds for any natural $k \in \N$ that $\mA^{2 k} = I$ and $\mA^{2 k + 1} = \mA$.

Let $\mA$ be an involution on a finite-dimensional Hilbert space.
If the involution $\mA$ is Hermitian, then it is also unitary.
For any complex $z \in \C$, the exponential $\dsp \exp\p{z \mA} = \sum_{k=0}^\infty \left(\frac{z^{2k}}{(2k)!} \mA^{2k} + \frac{z^{2k+1}}{(2k+1)!} \mA^{2k+1}\right)$ of $z \mA$ is given by $\exp\p{z \mA} = \left(\cosh z\right) I + \left(\sinh z\right) \mA$.
For any real $\theta \in \R$, the exponential $\dsp \exp\p{i\theta \mA} = \left(\cosh i \theta\right) I + \left(\sinh i \theta\right) \mA$ is given below.
\begin{align}
	\exp\p{i \theta \mA} = \left(\cos \theta\right) I + i \left(\sin \theta\right) \mA
\end{align}

